% !TEX root = main.tex
\section{Introduction}
\label{sec:introduction}

Time-series foundation models have become a practical alternative to models
trained separately for each forecasting application. Chronos-2
\citep{ansari2025chronos2}, TimesFM~2.5 \citep{das2025timesfm2}, Moirai~2.0
\citep{liu2025moirai2}, and TiRex \citep{auer2025tirex} can all be applied
without fitting their parameters to the target series. This makes them
attractive in retail, where a forecasting system may cover thousands of items
and where building and maintaining features for a supervised model can be
costly.

It is less clear whether this convenience comes with better forecasts. Model
papers and public benchmarks generally compare foundation models with one
another or with statistical methods. Retail applications, however, often rely
on global gradient-boosted trees with lags, calendar variables, prices, and
promotions. Published results rarely compare the two families on the same
series, forecast origins, accuracy measure, and information set. Differences
between benchmark tables may therefore reflect the evaluation design as much
as the forecasting method. This is especially important for intermittent
demand, for which metric aggregation and the treatment of zero sales can
change model rankings.

We study when a zero-shot foundation model is a useful alternative to a
conventional retail forecasting model. The analysis combines two forms of
evidence. First, we run five foundation models and several LightGBM and
statistical baselines on common 100-series samples from M5, Rohlik~v2, and
Favorita. Every model is evaluated on the same rolling forecast origins, and
prediction-level output is used for paired comparisons. We also vary the
LightGBM multi-step strategy and tuning budget, because a weak tree baseline
would exaggerate the apparent value of a foundation model. Second, we extract
retail tasks from fev-bench \citep{shchur2025fevbench} and GIFT-Eval
\citep{aksu2024gift}. These public results cover more tasks and models, but
they are analysed descriptively because their published tables do not contain
the prediction-level dependence information required for pooled statistical
inference.

The study addresses three questions:
\label{sec:research-questions}
\begin{enumerate}
  \item Do zero-shot foundation models improve retail demand forecasts relative
    to statistical and LightGBM baselines?
  \item Do the results vary with the accuracy measure, forecast horizon,
    demand pattern, and baseline specification?
  \item What evaluation procedure should a retailer use when choosing between
    the two model families?
\end{enumerate}

The matched experiments favour several foundation models, but not by a
constant margin. Chronos-2 has lower aggregate \WAPE{} than the best observed
local baseline in all nine dataset--horizon comparisons. TiRex also has a
lower point estimate in all nine, while results for the other models depend
more strongly on the dataset. The M5 comparison remains favourable to
Chronos-2 after richer LightGBM features, a larger tuning budget, and a more
intermittent sample are introduced. These findings apply to the tested panels
and feature sets; they are not comparisons with the best systems developed for
the M5 competition.

The public benchmarks give a more qualified picture. Across 498 matched
model--task--metric comparisons, foundation models have their clearest
advantage on Scaled Quantile Loss. Median differences are much smaller for
\WAPE{} and \MASE{}, and rankings vary across suites. The result is therefore
not a general ordering of model families. It shows instead that a useful
comparison requires a common evaluation panel, a suitable business metric,
and a credible conventional baseline.

The paper contributes matched retail experiments with prediction-level
uncertainty estimates, a reanalysis of two public benchmark suites, and an
evaluation procedure that practitioners can reproduce. It also documents how
LightGBM strategy, metric aggregation, covariates, and computational setting
affect the comparison. The following section places the models and benchmarks
in context. Sections~\ref{sec:methods}--\ref{sec:meta-regression} describe the
data and results, and Sections~\ref{sec:framework}--\ref{sec:conclusion}
discuss their practical and research implications.
